% ===============================================================
\section{Reasoning-aware Adaptive Routing Framework}\label{sec:routing_framework}
% ===============================================================

\subsection{Legal Knowledge Base}\label{subsec:legal_kb}

The vector index $\mathcal{I}_v$ contains $68{,}835$ embedded chunks
appropriate for direct questions whose answers are stated in legal text. The
graph $\mathcal{G}$ stores explicit typed relations among legal documents,
articles, and concepts. Tables~\ref{tab:graph_stats}
and~\ref{tab:graph_relations} report the statistics of the graph built from
the Pháp~điển codification, in which article nodes are promoted to
\textsc{LegalArticle} to enable article-level evidence retrieval.

\begin{table}[!htbp]
\centering
\caption{Vietnamese legal graph $\mathcal{G}$ statistics}\label{tab:graph_stats}
\begin{tabular}{lr}
\toprule
\textbf{Graph item} & \textbf{Count} \\
\midrule
Nodes $|\mathcal{V}|$ & 804{,}322 \\
Edges $|\mathcal{E}|$ & 1{,}968{,}621 \\
Legal documents & 4{,}598 \\
Legal articles (\textsc{LegalArticle}) & 68{,}570 \\
Semantic concepts & 392{,}592 \\
Other entities & 338{,}562 \\
\bottomrule
\end{tabular}
\end{table}

\begin{table}[!htbp]
\centering
\caption{Main relation types in the legal graph, ordered by count}\label{tab:graph_relations}
\begin{tabular}{lr}
\toprule
\textbf{Relation type} & \textbf{Count} \\
\midrule
\textsc{MENTIONS} & 1{,}033{,}341 \\
\textsc{HAS\_CONDITION} & 202{,}298 \\
\textsc{OBLIGATES} & 128{,}768 \\
\textsc{PART\_OF\_PROCEDURE} & 102{,}004 \\
\textsc{REGULATES} & 75{,}137 \\
\bottomrule
\end{tabular}
\end{table}

The graph is large and densely linked at the semantic level through relations
such as \textsc{MENTIONS} ($1{,}033{,}341$ edges) and \textsc{HAS\_CONDITION}
($202{,}298$), whereas explicit cross-document legal-effect relations such as
\textsc{AMENDS} ($5{,}944$ edges) and \textsc{REPLACES} ($4{,}189$) remain
comparatively sparse. This asymmetry directly motivates adaptive routing:
graph traversal is valuable for relation-heavy queries but unnecessarily
costly for simple article-level lookups.

\subsection{Reasoning-aware Feature Extraction}\label{subsec:features}

For each input $(q_t,H)$, the router first runs the history resolver
(Eq.~\ref{eq:history_resolver}) and then extracts a feature vector
$\mathbf{f}(q_t,H,r_H)\in\mathbb{R}^{n_f}$ that blends lexical signals with
reasoning-aware metrics, together with two auxiliary scores: the ambiguity
score $\phi_{\mathrm{amb}}(q_t,H,r_H)\in[0,1]$ and the multi-hop score
$\xi(q_t)\in\{0,1,2\}$.
Here $n_f=27$, comprising $16$ lexical surface features produced by a
Vietnamese legal feature extractor and $11$ reasoning-aware complexity
features produced by a query-complexity analyser; the complete list, in the
canonical vector order used for both training and inference, is given in
Appendix~\ref{app:repro}, Table~\ref{tab:features}. Using an explicit feature
vector rather than a hidden query embedding keeps Stage~1 both interpretable
and cheap to evaluate.
% TODO: re-export feature importances from the retrained 27-feature model
% (train_router.py writes them into the training report) and reinstate the
% top-4 numeric importances here; the previous values (0.3225/0.2131/...)
% belonged to the earlier model and were removed as stale.
The Stage~1 classifier relies most heavily on a small set of structural
features: the \textbf{legal reference count}, which tracks how many distinct
legal documents are mentioned; the \textbf{query length in words}, as
cross-document synthesis questions tend to be longer;
\textbf{cross-document signals}, reflecting explicit synthesis indicators;
and the \textbf{article reference count}. Additional reasoning
metrics---multi-hop connectors, sub-question count, procedural depth, and a
discrete complexity level---provide finer-grained separation for harder
cases, while question-type cues (\textit{gì}, \textit{khi nào},
\textit{bao nhiêu}) and legal-relation cues (\textit{sửa đổi},
\textit{thay thế}, \textit{có hiệu lực từ}) further inform the decision.

\paragraph{Multi-hop score.}
Let $\mathcal{K}_{\mathrm{rel}}$ be the fixed lexicon of legal-relation cue
phrases (e.g.\ \textit{sửa đổi}, \textit{thay thế}, \textit{bãi bỏ},
\textit{hướng dẫn}, \textit{căn cứ}), listed in full in
Appendix~\ref{app:repro}. The multi-hop score counts how many distinct cues
occur in the query, capped at two:
\begin{equation}
\xi(q_t)=\min\!\Bigl(2,\;\bigl|\{\,\kappa\in\mathcal{K}_{\mathrm{rel}} : \kappa\ \text{occurs in}\ q_t\,\}\bigr|\Bigr).
\label{eq:multihop_score}
\end{equation}
Thus $\xi(q_t)=0$ means no relation signal (likely a direct lookup),
$\xi(q_t)=1$ means a single relation cue, and $\xi(q_t)=2$ means two or more
co-occurring cues, indicating the strongest multi-hop demand.

\paragraph{Ambiguity score.}
The ambiguity score is produced by a deterministic detector that evaluates a
set of $m=11$ binary ambiguity indicators
$\{\psi_1(q_t,H,r_H),\ldots,\psi_m(q_t,H,r_H)\}$---for example, an
unresolved demonstrative or pronoun (\textit{đó}, \textit{này}, \textit{nó})
whose referent is not fixed by $r_H$, a generic legal need with no concrete
legal-object anchor, or an elliptical follow-up with no usable history---each
carrying a calibrated weight $w_i\in[0,1]$. The score is the strongest
triggered signal, gated by the history-resolution status:
\begin{equation}
\phi_{\mathrm{amb}}(q_t,H,r_H)=
\begin{cases}
0 & \text{if } r_H=\texttt{resolved},\\[3pt]
\displaystyle\max_{1\le i\le m}\; w_i\,\psi_i(q_t,H,r_H) & \text{otherwise},
\end{cases}
\label{eq:ambiguity_score}
\end{equation}
so that a referent recovered from history suppresses the ambiguity signal
regardless of surface cues. Taking the maximum rather than a sum reflects the
weight calibration: a single strong indicator (e.g.\ an unresolved pronoun
with empty history, $w=0.90$) is sufficient to exceed the trigger threshold
$\tau_{\mathrm{force}}=0.60$, whereas weak indicators ($w\le 0.30$) cannot
force verification on their own. The full indicator set and weights
$\{(\psi_i,w_i)\}$, including two graded indicator families whose scores
scale with match counts, are specified in Appendix~\ref{app:repro},
Table~\ref{tab:ambiguity_features}.

\paragraph{History resolver.}
A deterministic history resolver runs before ambiguity scoring,
\begin{equation}
r_H = \mathrm{Resolve}(q_t,H),
\label{eq:history_resolver}
\end{equation}
returning a referent-resolution status
\begin{equation}
r_H\in\{\texttt{resolved},\,\texttt{no\_history},\,\texttt{irrelevant\_history},\,\texttt{conflicting\_history},\,\texttt{not\_needed}\}.
\label{eq:history_states}
\end{equation}
If \textit{văn bản đó} (``that document'') has a unique document referent in
$H$, the resolved referent can unblock graph traversal; otherwise the same
surface form remains ambiguous. The status $r_H$ is used both as a Stage~1
input feature and as one of the Stage~2 trigger conditions; the resolution
rules are listed in Appendix~\ref{app:repro}.

\subsection{Stage~1: XGBoost Router}\label{subsec:stage1}

Stage~1 is an XGBoost multi-class classifier~\cite{chen2016xgboost} trained on
the $800$ routing-labelled training queries augmented with $60$ synthetic
\texttt{clarify} samples drawn from ambiguous-query templates, since the QA
dataset itself contains no gold \texttt{clarify} label:
\begin{equation}
(d_1,s_1) = f_1\bigl(\mathbf{f}(q_t,H,r_H)\bigr),
\label{eq:stage1}
\end{equation}
where $d_1\in\mathcal{D}=\{d_v,d_g,d_h,d_c\}$ is the initial route and
$s_1\in[0,1]$ is the confidence, taken as the maximum class probability of
the softmax output. Because the \texttt{clarify} class is only weakly
supervised by synthetic templates, a Stage~1 prediction $d_1=d_c$ is treated
as a proposal rather than a decision: whenever $d_1=d_c$, Stage~2 is invoked
to adjudicate, so clarification is never emitted on the strength of the
weakly supervised class alone.
% TODO: verify against two_stage_router.py that a Stage-1 clarify prediction
% unconditionally triggers Stage-2 (or is demoted when Stage-2 is off); if the
% actual mechanism differs, align this sentence and Algorithm 1 with the code. XGBoost is selected for three reasons: the
features are mostly discrete counts and binary indicators that
gradient-boosted trees handle without elaborate normalisation; inference over
$n_f$ features is orders of magnitude cheaper than one LLM pass, meeting the
latency-aware requirement of a stage that processes all traffic; and tree
models expose feature importances that support interpretability, which is
valuable in the legal domain. Empirically, the classifier forward pass takes
$\approx\!8.9$~ms per query on CPU (the full Stage~1 path, including feature
extraction and history resolution, is $\approx\!33.7$~ms; see
Section~\ref{subsec:rq5}), versus $\approx\!53.3$~ms for a
PhoBERT~\cite{nguyen2020phobert} classifier on GPU, making Stage~1 suitable as
a real-time gateway and motivating a future three-level cascade
XGBoost~$\rightarrow$~PhoBERT~$\rightarrow$~LLM following
FrugalGPT~\cite{chen2023frugalgpt}.

\subsection{Stage~2: Selective LLM Verifier}\label{subsec:stage2}

Stage~2 is an LLM verifier invoked selectively. Three binary trigger
indicators, each taking values in $\{0,1\}$, are defined as
\begin{align}
\gamma_{\mathrm{conf}} &= \mathbb{I}[\,s_1 < \tau_{\mathrm{conf}}\,],
    \label{eq:gamma_conf}\\
\gamma_{\mathrm{amb}}  &= \mathbb{I}[\,\phi_{\mathrm{amb}}(q_t,H,r_H) \geq \tau_{\mathrm{force}}\,],
    \label{eq:gamma_amb}\\
\gamma_{\mathrm{hop}}  &= \mathbb{I}[\,\xi(q_t) \geq 2\,],
    \label{eq:gamma_hop}
\end{align}
where $\mathbb{I}[\cdot]$ is the indicator function (equal to $1$ when its
argument is true and $0$ otherwise). Thus $\gamma_{\mathrm{conf}}$ fires when
Stage~1 confidence is below $\tau_{\mathrm{conf}}$, $\gamma_{\mathrm{amb}}$
fires when the ambiguity score reaches the forcing threshold
$\tau_{\mathrm{force}}$, and $\gamma_{\mathrm{hop}}$ fires at the strongest
multi-hop level ($\xi=2$). With the convention $a\vee b=\max(a,b)$ on
$\{0,1\}$, the composite trigger is the logical disjunction of the three
indicators together with the conflicting-history condition,
\begin{equation}
\Gamma(q_t,H,r_H) = \gamma_{\mathrm{conf}}\vee\gamma_{\mathrm{amb}}\vee\gamma_{\mathrm{hop}}\vee\mathbb{I}[\,r_H=\texttt{conflicting\_history}\,]\;\in\;\{0,1\},
\label{eq:trigger_function}
\end{equation}
so that $r_H$ enters the trigger explicitly.

When $\Gamma=0$, Stage~1 is trusted and $d^*=d_1$. When $\Gamma=1$, the
verifier is queried and returns a proposed route $\tilde{d}\in\mathcal{D}$
together with, when $\tilde{d}=d_c$, a clarify confidence
$p_c\in[0,1]$ (the model's estimated probability that the query cannot be
safely grounded as posed). A clarification is emitted only if the verifier is
sufficiently confident; otherwise the proposed retrieval route is kept. The
clarify-override threshold $\tau_{\mathrm{clar}}$ therefore acts on the
verifier output, and the final route is
\begin{equation}
d^* =
\begin{cases}
d_1 & \text{if } \Gamma(q_t,H,r_H)=0,\\[3pt]
d_c & \text{if } \Gamma=1,\ \tilde{d}=d_c,\ \text{and } p_c\ge\tau_{\mathrm{clar}},\\[3pt]
\tilde{d} & \text{otherwise } (\Gamma=1,\ \tilde{d}\in\{d_v,d_g,d_h\}\ \text{or } p_c<\tau_{\mathrm{clar}}),
\end{cases}
\label{eq:final_route}
\end{equation}
where in the last case a below-threshold clarify proposal is demoted to the
verifier's best retrieval route. Default thresholds are
$\tau_{\mathrm{conf}}=0.50$, $\tau_{\mathrm{force}}=0.60$, and
$\tau_{\mathrm{clar}}=0.80$; their effect is analysed in
Section~\ref{subsec:rq4}. In words, Stage~2 is invoked on low-confidence
Stage~1 decisions, on high ambiguity risk, on strong multi-hop signals, or on
conflicting history; a \texttt{resolved} history unblocks retrieval and
suppresses false clarification, whereas missing, irrelevant, or conflicting
history forces verification or clarification.

Algorithm~\ref{alg:router} summarises the procedure. The deployed
implementation additionally uses two heuristics: a dense fast path that skips
Stage~2 for high-confidence, low-ambiguity dense queries, and a
reasoning-signal check that escalates queries with strong legal features even
at moderate Stage~1 confidence.

\begin{algorithm}[!htbp]
\caption{Two-stage adaptive router}\label{alg:router}
\begin{algorithmic}[1]
\Require query $q_t$, history $H$, thresholds $\tau_{\mathrm{conf}},\tau_{\mathrm{force}},\tau_{\mathrm{clar}}$
\Ensure final route $d^*\in\mathcal{D}$
\State $r_H \gets \mathrm{Resolve}(q_t,H)$
\State $(\mathbf{f},\phi_{\mathrm{amb}},\xi) \gets \mathrm{ExtractFeatures}(q_t,H,r_H)$
\State $(d_1,s_1) \gets f_1(\mathbf{f})$
\State $\gamma_{\mathrm{conf}}\gets\mathbb{I}[s_1<\tau_{\mathrm{conf}}]$;\ \ $\gamma_{\mathrm{amb}}\gets\mathbb{I}[\phi_{\mathrm{amb}}\ge\tau_{\mathrm{force}}]$;\ \ $\gamma_{\mathrm{hop}}\gets\mathbb{I}[\xi\ge 2]$
\State $\Gamma \gets \gamma_{\mathrm{conf}}\vee\gamma_{\mathrm{amb}}\vee\gamma_{\mathrm{hop}}\vee\mathbb{I}[r_H=\texttt{conflicting\_history}]$
\If{$\Gamma=0$}
  \State \Return $d_1$
\EndIf
\State $(\tilde{d},p_c) \gets f_2(q_t,H,d_1,s_1,\mathbf{f})$
\If{$\tilde{d}=d_c$ \textbf{and} $p_c\ge\tau_{\mathrm{clar}}$}
  \State \Return $d_c$
\Else
  \State \Return $\tilde{d}$ \Comment{demote below-threshold clarify to best retrieval route}
\EndIf
\end{algorithmic}
\end{algorithm}

\noindent\textbf{Running example.}
For $q_1=$ \textit{``Điều kiện để được cấp giấy phép xây dựng là gì?''}
(``What are the conditions for a building permit?''), $\xi(q_1)=0$ and Stage~1
predicts $d_v$ with high confidence; Stage~2 is bypassed. For $q_2=$
\textit{``Nghị định nào sửa đổi quy định về xử phạt vi phạm hành chính trong
lĩnh vực đất đai?''} (``Which decree amends the administrative-penalty rules
on land?''), the phrase \textit{sửa đổi} (``amend'') raises relational
signals, so the router selects $d_g$ or $d_h$. For $q_3=$ \textit{``Văn bản đó
còn hiệu lực không?''} (``Is that document still in effect?'') with
$H=\emptyset$, the referent is unresolved, $\phi_{\mathrm{amb}}$ is high,
$\Gamma=1$, and the verifier returns $\tilde{d}=d_c$ with $p_c\ge\tau_{\mathrm{clar}}$.

\subsection{Route Execution and Fallback Policy}\label{subsec:route_execution}

For $d_v$, the system retrieves top-$k$ chunks from $\mathcal{I}_v$. For
$d_g$, it links query entities to graph nodes and traverses $\mathcal{G}$. For
$d_h$, it constructs $E_{d^*}=E_v\cup E_g$. For $d_c$, it returns a
clarification question without retrieval. A defensive fallback policy applies
when primary retrieval is uninformative: a failed graph traversal or hybrid
synthesis falls back to dense retrieval; an uninformative dense lookup expands
$k$ once and retries; and a clarify route, if unanswered, keeps the
conversation in a waiting state rather than guessing the referent.
